\chapter{Individual application regression tests}\label{ch:application-regression}
\section{The legacy application procedure and its authorities}
AMPS also retains application regressions generated from
\path{MakefileTest/Table} into \path{Makefile.test}. These tests configure an
application, build a native executable, create a run directory, evolve it and
compare selected output with maintained references. They complement the
modern core/application registry of Chapter~\ref{ch:tests}; they are not
automatically imported into that Python runner.

The edition snapshot has 153 Table records, describing 150 distinct targets.
Three mover names have alternative machine-specific declarations. The
generated \code{HOSTNAME=passion} file contains 100 application owners;
50 other declared targets are absent from this host snapshot.
An omitted record is not an executed SKIP. A generated target is not
necessarily enabled for the current compiler, provisioned with data, or
physically qualified.

Appendix~\ref{ch:legacy-catalog} describes all 100 generated owners
individually, with source line numbers, named purpose, actual comparison type,
explicit tolerance overrides and detected prerequisites. It separately lists
all 50 omitted names and their machine selectors. The companion
\path{doc/legacy_test_catalog.json} retains exact phase locations and both
source hashes. These counts are an editorial inventory, not a new test run.

\section{How Table records control generation}
The maintained record is delimited by \code{#>} and \code{<#}.
\begin{longtable}{L{0.22\textwidth}L{0.68\textwidth}}
\caption{Application-regression declaration fields.}\\
\toprule Field & Meaning\\\midrule\endfirsthead
\toprule Field & Meaning\\\midrule\endhead
\code{Name} & Application/test path. The generator removes the
\code{test/} prefix from the target name. Example:
\code{Name=test/fast-wave} produces \code{test_fast-wave}.\\
\code{Time} & Declared approximate runtime used by the historical scheduling
tools; not a measured duration or an enforced physical acceptance limit.\\
\code{Type} & \code{amps} or \code{swmf}; distinguishes a standalone case
from a custom framework-coupled campaign.\\
\code{Keys} & Configuration options, frequently selecting SPICE, background
data, OpenMP or vectorization. Machine-qualified alternatives are possible.\\
\code{Outs} & Declared output basenames for the default comparison.
When omitted, the default basename is \code{test_NAME}.\\
\code{Ref} & Compiler/machine-dependent reference suffix, for example a
number in brackets inserted before the rank-specific reference suffix.\\
\code{ExeptionCode} & Historical spelling retained by the generator.
A compiler-dependent SKIP can suppress aggregate execution. It is not a
physical PASS and a direct target has its own actual behavior.\\
\path{CustomRefSolutionPaths} & Site/reference locations used by custom
recipes where defined; inspect the effective recipe for actual use.\\
\code{Compile}, \code{Rundir}, \code{Run}, \code{Check} & Optional replacement
recipes. Lines beginning with \code{>>>} acquire a recipe TAB; \code{<<<}
preserves make-level text such as conditionals.\\
\bottomrule
\end{longtable}

A missing machine selector applies everywhere. An inclusion selector such
as \code{@@(0+passion+pleiades)} names permitted machines; an exclusion
selector such as \code{@@(-pleiades)} excludes named machines.
\code{@@(0+)} names no permitted host. Selection and precedence are implemented
by the production generator; do not infer a case's effective options simply
by concatenating every matching-looking Table line.

The current \path{Config.pl} installation and \code{-test-blocks=N} paths
compile/run \path{utility/TestScripts/table_parser.cpp}. The older Perl
\path{utility/TestScripts/BuildTest.pl} and template files are historical
implementation references and differ in some behavior. Installation normally
generates two execution blocks; changing the count regenerates the file.
Do not edit only the generated Makefile and expect a change to survive the
next generation. Keep the source declaration and effective recipe consistent.

\section{Four phases and an individual invocation}
An owner normally provides:
\begin{lstlisting}
test_NAME
test_NAME_compile
test_NAME_rundir
test_NAME_run
test_NAME_check
\end{lstlisting}
The aggregate target calls those phases in order. The compile phase selects
the application and libraries. Rundir usually recreates
\path{run_test_NAME} and moves the executable into it. Run advances that
frozen executable. Check evaluates the declared files and writes diff/log
artifacts. Each phase needs its predecessor's exact state; an old run
directory is not proof that the current binary ran.

For the two currently preflight-adapted regression examples:
\begin{lstlisting}
# From the idle AMPS root; choose a launcher supported by the host.
make test_Moon MPIRUN='mpiexec -n 4'
make test_fast-wave MPIRUN='mpiexec -n 4'
\end{lstlisting}
Their compile recipes invoke
\path{src/test/native/prepare_legacy_native.py}, then the established
configuration and top-level -j16 build. Some other custom recipes hard-code
MPI rank counts or run nested SWMF campaigns, so overriding \code{MPIRUN}
does not universally change their resources. Read the individual run recipe.

The fast-wave recipe alone adds a GNU host C/C++ \code{-O1} fallback if the
site's effective flags specify no optimization level. An explicit level,
NVCC or another compiler family keeps its existing policy. Disable the
fallback with \code{FAST_WAVE_CXXFLAGS=}; site files and Fortran flags are
preserved. On passion the complete clean four-rank command improved from
119.04 to 87.08 seconds, with all six sampled species files byte-identical.
This dated measurement retains the original comparison bounds; it is not a
portable speed threshold. See \path{test/srcFastWave/README.md} for the exact
timing method, controls and evidence.

The generated \code{test_help} target is intended to list these names.
In this edition, executing \code{make test_help} returns a shell quoting error:
some generated echo lines lack their closing quote. Use the complete local
catalog and actual target definitions until that separate generator defect
is repaired. The manual documents the defect; it does not report the help
command as passing.

\section{Application and numerical families}
\begin{longtable}{L{0.29\textwidth}L{0.61\textwidth}}
\caption{The 100 currently generated application owners.}\\
\toprule Family / owners & What the fixtures exercise\\\midrule\endfirsthead
\toprule Family / owners & What the fixtures exercise\\\midrule\endhead
Planetary/comet/dust: 13 & Saturn MOP; Moon baseline/cut-cell/restart;
Enceladus and Mercury; CG gas/dust and tetrahedral variants; model dust;
two Europa cases and Venus hot oxygen.\\
Heliospheric/coupled: 29 & Parker equation, focused transport, Boris and
field-line transport; 3-D prescribed or SWMF IMF; SC/IH field-line/node
attachment and diffusion variants; ARMS background; BATS-R-US PIC, Europa,
outer-heliosphere and dynamic-flux-rope couplings.\\
Fields/waves/solvers: 22 & Linear solver and OpenMP variant, GMRES,
uniform/nonuniform periodic boundaries, current-on/off light waves,
reconnection, eleven fast-wave selector/vector/thread/rebalancing variants
and constant-field boundary behavior.\\
Trajectory/mover: 12 & Individual Lapenta2017, Markidis2010, Boris,
relativistic Boris and guiding-centre paths; CCMC; corresponding mover
applications and relativistic guiding-centre fixture.\\
Numerical/process: 10 & Four legacy GoogleTest applications, NTC/MF
collision distributions, two chemistry references, AMR interpolation through
two adapters and CDAP diagnostic.\\
Geospace: 14 & Gridless/mesh-backed and trajectory cutoff rigidity,
vertical/hybrid variants, analytic dipole/T96 forward products, gridless
directional density/flux fixtures and live SWMF/PT products.\\
\bottomrule
\end{longtable}
The legacy Enceladus moment comparison is different from the modern
10000-step E3 source/observation owner. The legacy Moon baseline is different
from terrain/science qualification. A modern global mapping/audit that
mentions a legacy directory does not execute every Table application:
\code{CORE-LEGACY-REGISTRY-AUDIT-001} is a separate 13-directory governance
audit, not the 100-case regression procedure.

\section{References, tolerances and comparison artifacts}
The default recipe compares
\path{run_test_NAME/PT/plots/OUT.dat} with a maintained file under
\path{output/test_NAME}, using the selected reference suffix and the actual
rank-specific \code{ref_npN} convention. The historical \path{output} can be
a link to a shared reference tree. Preserve those references and record their
identity before a run.

Custom checks can reduce rows, remove nondeterministic columns, compare
trajectories/surfaces/volumes separately or import SWMF's own diff artifacts.
The current catalog detects DiffNum in 93 owners, plain diff in three,
and four custom checks with neither direct invocation. Those four include
SWMF diff imports and the legacy unit-test stage/log check.

\path{share/Scripts/DiffNum.pl} uses its own default absolute/relative
criteria unless a recipe overrides them. The appendix reproduces explicit
\code{-a=} and \code{-r=} settings; it does not silently substitute the
modern core suite's tolerances. Plain diff is a text equality check.
Compiler- and rank-specific references must match the actual executable and
layout. Do not assume stochastic particle records are rank-identical.

An empty final diff means agreement only if every prerequisite comparison
actually ran and required files exist. Check each expected sub-diff and the
final combined result, as well as process completion. Some default
multi-output paths redirect repeatedly to one diff; a last empty comparison
must not hide an earlier failure. Some custom recipes conditionally copy a
framework result. Missing upstream evidence is not a successful comparison.

The base \code{test_all} recipe suppresses some child make failures and
clears old diff files. The root Makefile also adds a modern core-nightly
prerequisite to \code{test_all}, so the effective target can combine both
procedures. Use the explicit Python selection for the modern registry and an
explicit reviewed individual legacy target for a bounded regression.
Never infer all tests passed solely from the outer make return code.

\section{Dependencies, preflight and concurrency}
Site paths in \code{Keys} can require CSPICE/kernels, background archives,
Kameleon, GoogleTest or a separately installed SWMF tree. Vector/thread
variants require matching compiler, CPU and OpenMP settings. The catalog
flags detected prerequisites; it is not a complete environment solver.
Some recipes also copy files produced by another test: Moon-restart explicitly
requires particle restart files from the matching Moon baseline.

Many historical compile recipes remove \path{srcTemp}, use -j8, or switch
applications without the current required root/idle/build-removal preflight.
They must be reconciled before execution under the present policy.
The shared native preflight helper is explicitly present in the current Moon
and fast-wave recipes; its presence does not qualify every other target.
Twenty-four generated owners contain Git instructions in custom recipes.
The current Git-read-only policy does not permit executing those commands;
documentation of them is not permission to clone, switch, pull or update.

Historical execution blocks are not the Python runner's \code{--streams}
abstraction. In particular, \code{test_compile} and
\code{test_run_threadN} do not establish a safe cross-stream native-build
lock. Do not run two legacy configuring/compiling targets concurrently in
one checkout. A future global adapter must serialize mutable build work,
freeze executables/inputs, retain per-case logs and faithfully import all
FAIL/SKIP/ERROR states before combining reports.

\section{Maintaining this documentation}
After changing either authoritative file, refresh the local snapshot:
\begin{lstlisting}
python3 -B doc/update_legacy_test_catalog.py
python3 -B doc/update_legacy_test_catalog.py --check
make -C doc check
make PDF
\end{lstlisting}
This editorial generator only reads declarations and recipes; it never
regenerates \path{Makefile.test} or executes its commands. Add a reviewed
purpose description for any new generated owner. Preserve all declarations,
including machine-omitted entries, and never turn catalog membership into
a reported runtime PASS.
\source{MakefileTest/Table; Makefile.test; utility/TestScripts/table_parser.cpp}
